\chapter{Emacs\index{Emacs} Customization}

This appendix provides some tips and techniques for customizing Emacs to work more effectively with Common Lisp\index{Common Lisp}.

\section{Key Bindings}

The default Emacs key bindings use Control and Meta extensively. On most keyboards:

\begin{itemize}
\item \textbf{Control} is the key labeled ``Ctrl'' or ``Control''
\item \textbf{Meta} is typically the key labeled ``Alt'' or (on some systems) the Windows key
\end{itemize}

\subsection{Remapping Caps Lock}

Many Emacs users find it convenient to remap the Caps Lock key to function as an additional
Control key, since Control is used so frequently. This can significantly reduce hand strain
for heavy Emacs users.

\subsubsection{On Linux/X11}

Add the following to your \texttt{.xinitrc} or \texttt{.xsession} file:

\begin{verbatim}
setxkbmap -option ctrl:nocaps
\end{verbatim}

Or use \texttt{xmodmap} with a file \texttt{.xmodmaprc} containing:

\begin{verbatim}
clear Lock
keycode 66 = Control_L
add Control = Control_L
\end{verbatim}

Then run:

\begin{verbatim}
xmodmap ~/.xmodmaprc
\end{verbatim}

\subsubsection{On Windows}

Several utilities are available to remap keys on Windows, such as:

\begin{itemize}
\item SharpKeys (free, GUI-based)
\item AutoHotkey (scriptable)
\item Windows Registry editing (for advanced users)
\end{itemize}

\subsubsection{On macOS}

Go to System Preferences $\rightarrow$ Keyboard $\rightarrow$ Modifier Keys, and remap Caps Lock to
Control.

\section{Essential Emacs Commands}

Here are some essential Emacs commands that are particularly useful when working with
Lisp:

\subsection{Navigation}

\begin{itemize}
\item \texttt{C-a} --- Move to beginning of line
\item \texttt{C-e} --- Move to end of line
\item \texttt{M-<} --- Move to beginning of buffer
\item \texttt{M->} --- Move to end of buffer
\item \texttt{C-M-f} --- Move forward by s-expression
\item \texttt{C-M-b} --- Move backward by s-expression
\item \texttt{C-M-d} --- Move down into s-expression
\item \texttt{C-M-u} --- Move up out of s-expression
\end{itemize}

\subsection{Editing}

\begin{itemize}
\item \texttt{C-k} --- Kill (cut) to end of line
\item \texttt{M-w} --- Copy region
\item \texttt{C-w} --- Kill (cut) region
\item \texttt{C-y} --- Yank (paste)
\item \texttt{C-/} or \texttt{C-\_} --- Undo
\item \texttt{C-M-k} --- Kill s-expression forward
\item \texttt{M-C-k} --- (alternate) Kill s-expression forward
\item \texttt{C-M-t} --- Transpose s-expressions
\end{itemize}

\subsection{Buffers and Windows}

\begin{itemize}
\item \texttt{C-x b} --- Switch to buffer
\item \texttt{C-x k} --- Kill (close) buffer
\item \texttt{C-x C-f} --- Find (open) file
\item \texttt{C-x C-s} --- Save file
\item \texttt{C-x C-w} --- Write file (save as)
\item \texttt{C-x 2} --- Split window horizontally
\item \texttt{C-x 3} --- Split window vertically
\item \texttt{C-x 1} --- Delete other windows
\item \texttt{C-x 0} --- Delete current window
\item \texttt{C-x o} --- Switch to other window
\end{itemize}

\subsection{Lisp-Specific Commands}

When using the Emacs-Lisp interface (as described in Chapter 2):

\begin{itemize}
\item \texttt{C-M-x} --- Compile and load the current top-level form
\item \texttt{C-c C-b} --- Compile and load the current buffer
\item \texttt{M-A} --- Show argument list for function at point
\item \texttt{C-c .} --- Find definition of symbol at point
\item \texttt{M-,} --- Return from definition lookup
\item \texttt{C-c C-k} --- Compile and load file
\item \texttt{C-c C-z} --- Switch to Lisp listener
\end{itemize}

\section{Customizing .emacs}

The \texttt{.emacs} file in your home directory is executed when Emacs starts up. This file can
contain Emacs Lisp code to customize your Emacs environment.

\subsection{Basic .emacs Example}

Here is a simple \texttt{.emacs} file with some useful settings:

\begin{verbatim}
;; Load the Allegro CL\index{Allegro CL} - Emacs interface
(setq load-path (cons "/usr/local/acl62/eli" load-path))
(load "fi-site-init")

;; Don't show the startup message
(setq inhibit-startup-message t)

;; Show line and column numbers
(setq line-number-mode t)
(setq column-number-mode t)

;; Enable syntax highlighting
(global-font-lock-mode t)

;; Show matching parentheses
(show-paren-mode t)

;; Electric pair mode (auto-close parens)
(electric-pair-mode t)

;; Auto-revert buffers when files change on disk
(global-auto-revert-mode t)

;; Use spaces instead of tabs
(setq-default indent-tabs-mode nil)

;; Set tab width to 2 spaces
(setq-default tab-width 2)
\end{verbatim}

\subsection{Keyboard Macro for Common Lisp Session}

You can define custom key bindings to make frequent operations more convenient:

\begin{verbatim}
;; Quick key to switch to *common-lisp* buffer
(global-set-key (kbd "C-c l")
                (lambda ()
                  (interactive)
                  (switch-to-buffer "*common-lisp*")))

;; Quick key to start Common Lisp
(global-set-key (kbd "C-c C-l") 'fi:common-lisp)
\end{verbatim}

\subsection{ParEdit Mode}

ParEdit is a minor mode for Emacs that provides structured editing of S-expressions. It
ensures that parentheses always remain balanced and provides commands for manipulating
S-expressions structurally. Many experienced Lisp programmers find ParEdit indispensable.

ParEdit can be installed via MELPA (see Section \ref{sec:melpa}).

\subsection{Rainbow Delimiters}

Rainbow Delimiters mode colorizes nested parentheses with different colors, making it easier
to see the structure of deeply nested code. This is particularly helpful when working with
Lisp.

To use Rainbow Delimiters, install it via MELPA and add to your \texttt{.emacs}:

\begin{verbatim}
(require 'rainbow-delimiters)
(add-hook 'lisp-mode-hook 'rainbow-delimiters-mode)
(add-hook 'emacs-lisp-mode-hook 'rainbow-delimiters-mode)
\end{verbatim}

\section{Package Management: MELPA}
\label{sec:melpa}

MELPA (Milkypostman's Emacs Lisp Package Archive) is a repository of Emacs packages
that can be easily installed. To enable MELPA, add the following to your \texttt{.emacs}:

\begin{verbatim}
(require 'package)
(add-to-list 'package-archives
             '("melpa" . "http://melpa.org/packages/") t)
(package-initialize)
\end{verbatim}

After restarting Emacs, you can use \texttt{M-x package-list-packages} to browse and install
packages.

\section{SLIME\index{SLIME}: An Alternative to the Emacs-Lisp Interface}

SLIME (Superior Lisp Interaction Mode for Emacs) is an alternative to the Franz Inc.
Emacs-Lisp interface. SLIME is a popular open-source project that works with many Common Lisp implementations, including SBCL, CCL, and Allegro CL.

SLIME can be installed via MELPA or Quicklisp. It provides features such as:

\begin{itemize}
\item Interactive debugging with a graphical debugger
\item Code completion
\item Cross-referencing
\item Inspector for examining objects
\item Profiling and tracing support
\end{itemize}

For more information, visit:

\begin{center}
\texttt{http://common-lisp.net/project/slime/}
\end{center}

\section{Color Themes}

Emacs supports color themes to customize the appearance of your editing environment. Many
themes are available through MELPA, or you can create your own.

Popular themes include:

\begin{itemize}
\item solarized (light and dark versions)
\item zenburn
\item monokai
\item tomorrow (light and dark versions)
\end{itemize}

To try a theme, install it via MELPA and use \texttt{M-x load-theme}.

\section{Other Useful Packages}

\subsection{Company Mode}

Company (Complete Anything) provides text completion. It works well with SLIME for
Common Lisp code completion.

\subsection{Magit}

Magit is a sophisticated interface to Git version control within Emacs. If you're using Git
for your Lisp projects, Magit is highly recommended.

\subsection{Projectile}

Projectile provides project management features, making it easy to navigate between files in
a project, search within a project, and more.

\section{Learning More About Emacs}

Emacs has extensive built-in documentation:

\begin{itemize}
\item \texttt{C-h t} --- Start the Emacs tutorial
\item \texttt{C-h k} --- Describe key binding
\item \texttt{C-h f} --- Describe function
\item \texttt{C-h v} --- Describe variable
\item \texttt{C-h m} --- Describe current major and minor modes
\item \texttt{C-h i} --- Open Info documentation browser
\end{itemize}

For online resources, the Emacs Wiki (\texttt{http://www.emacswiki.org/}) contains a wealth
of information, tips, and customization examples.
